%\begin{abstract}
% Foreword
%     Context:
%     Need:
%     Task (what the authors did):
%     Object (what the document does):
% Summary
%     Findings:
%     Conclusion:
%     Perspectives

%Este documento provee de la información e instrucción necesarias para preparar el informe de la tarea para el seminario de Modelación computacional de materiales estructurales. El presente resumen no debe contener más de 300 palabras y debe estar estructurado por un prefacio y un sumario. En específico, el prefacio debe contener la motivación detrás del trabajo:

%\begin{itemize}
%\item Contexto del trabajo.
%\item Necesidad de simular virtualmente los ensayos.
%\item Tarea realizada por los autores (simulación de un ensayo específico para un material).
%\item Objeto del informe; es decir lo que el documento contiene.
%\end{itemize}
%
% sumario, en tanto, debe contener el aporte de los autores:
%
%\begin{itemize}
%\item Resultados relevantes.
%\item Conclusionesa partir de estos resultados.
%\item Opcionalmente, perspectivas que se puedan avizorar a partir del trabajo.
%\end{itemize}

%\paragraph{Nota:} El resumen no debe contener saltos de líneas. Todo lo anteriormente descrito debe estar contenido en un solo párrafo. 
%\end{abstract}



%\section{Introducción}

%Un documento debe estar organizado en tres partes:

%\begin{enumerate}
%\item Introducción 
%\item Cuerpo del trabajo.
%\item Conclusiones.
%\end{enumerate}

%Recordar que cada párrafo es una idea. Tratar de ser claro, riguroso y conciso. 


\section{Problema y objetivo}
El objetivo de este informe técnico es comparar experimentalmente la deflexión medida en el centro de la luz de una viga simplemente apoyada con las predicciones analíticas obtenidas mediante la teoría de Euler-Bernoulli, evaluando la respuesta lineal elástica del elemento bajo un régimen de cargas puntuales crecientes de 0 a 40 kN.

\section{Datos y supuestos}
Se analiza una viga de sección transversal rectangular ($b = 200\text{ mm}$, $h = 400\text{ mm}$), luz $L = 4000\text{ mm}$ y módulo de elasticidad $E = 25\text{ GPa}$. Se asumen las hipótesis de Euler-Bernoulli (secciones planas permanecen planas y deformaciones pequeñas). Las cargas centradas ($P$) varían de 0 a 40 kN en escalones de 5 kN.

\begin{figure}[htbp]
    \centering
    \includegraphics[width=0.65\textwidth]{esquema_viga.png}
    \caption{Esquema de la viga simplemente apoyada y sección transversal.}
    \label{fig:esquema}
\end{figure}

\section{Método}

Las conversiones de unidades utilizadas para unificar el cálculo en kN y mm corresponden a:
\[
1000\text{ mm} = 1\text{ m} \quad \text{y} \quad 1\text{ GPa} = 1\text{ kN/mm}^2
\]

El segundo momento de área ($I$) de la sección transversal rectangular se calcula como:
\begin{equation}
I = \frac{b \cdot h^3}{12}
\label{eq:inercia}
\end{equation}

La deflexión teórica en el centro de la luz ($\delta_{teo}$) para cada nivel de carga responde a la expresión:
\begin{equation}
\delta_{teo} = \frac{P \cdot L^3}{48 \cdot E \cdot I}
\label{eq:deflexion}
\end{equation}

Como indicador de comportamiento elástico lineal, la pendiente teórica o razón entre deflexión y carga ($\delta/P$) se determina mediante:
\begin{equation}
\frac{\delta}{P} = \frac{L^3}{48 \cdot E \cdot I}
\label{eq:flexibilidad}
\end{equation}

Para cargas $P \neq 0$, se evalúa la diferencia relativa porcentual:
\begin{equation}
\text{Diferencia relativa (\%)} = \left| \frac{\delta_{medida} - \delta_{teo}}{\delta_{teo}} \right| \times 100\%
\label{eq:diferencia}
\end{equation}

\section{Resultados y verificación}

Sustituyendo las dimensiones geométricas con las conversiones de unidades correspondientes en (1), se obtiene:
\[
I = 1{,}067 \times 10^9\text{ mm}^4
\]

Evaluando la Ecuación (3), la pendiente teórica deflexión-carga resulta:
\[
\frac{\delta}{P}  = 0{,}050\text{ mm/kN}
\]

En el Cuadro 1  se muestra detalladamente la comparación entre los valores medidos y teóricos.


\begin{table}[htbp]
    \centering
    \caption{Comparación teórico-experimental de deflexiones, relación $\delta/P$ y diferencia relativa.}
    \label{tab:resultados}
    \vspace{0.2cm}
    \begin{tabular}{ccccc}
        \toprule
        $P$ (kN) & $\delta_{medida}$ (mm) & $\delta_{teorica}$ (mm) & $\delta/P$ (mm/kN) & Diferencia relativa \\
        \midrule
       0  & 0,00 & 0,00 & --    & -- \\
        5  & 0,26 & 0,25 & 0,05 & 4\% \\
        10 & 0,50 & 0,50 & 0,05 & 0\% \\
        15 & 0,76 & 0,75 & 0,05 & 1\% \\
        20 & 1,01 & 1,00 & 0,05 & 1\% \\
        25 & 1,28 & 1,25 & 0,05 & 2\% \\
        30 & 1,54 & 1,50 & 0,05 & 3\% \\
        35 & 1,77 & 1,75 & 0,05 & 1\% \\
        40 & 2,05 & 2,00 & 0,05 & 2\% \\
        \bottomrule
    \end{tabular}
\end{table}

Los resultados confirman una concordancia destacable: la diferencia relativa entre la deflexión medida y la teórica no supera el 4\% en ningún escalón de carga. Asimismo, la constancia de la razón $\delta/P \approx 0{,}05\text{ mm/kN}$. 

\begin{figure}[htbp]
    \centering
    \includegraphics[width=0.60\textwidth]{carga_deflexion.png}
    \caption{Curva Carga vs. Deflexión: comparación teórico-experimental.}
    \label{fig:curva_deflexion}
\end{figure}
El gráfico muestra claramente un comportamiento lineal ascendente, lo que  verifica numericamente, en base a las cargas aplicadas, que el elemento opera dentro del régimen elástico lineal, manteniendo su rigidez sin signos de plastificación.

Como comprobación independiente para un caso particular de carga ($P = 20\text{ kN}$), al sustituir directamente en (2) se obtiene:
\[
\delta_{teo}(20\text{ kN}) = \frac{20\text{ kN} \cdot (4000\text{ mm})^3}{48 \cdot 25\text{ kN/mm}^2 \cdot (1{,}067 \times 10^9\text{ mm}^4)} = 1{,}00\text{ mm}
\]
Lo cual coincide exactamente con el valor teórico tabulado.

\section{Interpretación y conclusiones}

\textbf{Interpretación y conclusiones:} La teoría de Euler-Bernoulli predice con alta precisión las deflexiones medidas, registrando diferencias relativas inferiores al 4\%. La constante $\delta/P = 0{,}05\text{ mm/kN}$ confirma que la viga opera en régimen elástico lineal en todo el rango evaluado (0–40 kN).

\vspace{0.5cm}

\textbf{Limitaciones:} El modelo omite la deformación por esfuerzo cortante y asume apoyos ideales. No obstante, dada la esbeltez de la viga ($L/h = 10$), no altera la validez del cálculo.

\section{Declaración de uso de IA}
Se utilizó inteligencia artificial como herramienta de apoyo para realizar el paper en código en LaTeX y revisión estilística. Los cálculos, gráficos y análisis técnico fueron elaborados por el autor.

\begin{thebibliography}{9}
\bibitem{bachmann2026}
J. Bachmann, \textit{Diapositivas y apuntes de clases: Análisis Estructural}, Universidad de Santiago de Chile, 2025.

\bibitem{hibbeler2016}
R. C. Hibbeler, \textit{Estática}, 14.ª ed., Pearson Educación, 2016.
\end{thebibliography}


\end{document}